% \newpage
% \section{Trabajo Extra Voluntario}
% \label{sec:trabajo_extra}

% Esta sección recoge las ampliaciones voluntarias sobre la base de la práctica 1 y
% su adaptación al problema del portfolio de esta segunda práctica. El punto de
% partida heredado es el bloque formado por Greedy, Random y la Búsqueda Local
% primer mejor, que ya permite contrastar exploración pura, construcción heurística
% y explotación local. Sobre esa base se han incorporado variantes que refinan la
% búsqueda local, cambian el operador de mutación y prueban una familia distinta de
% algoritmos evolutivos.

% \subsection{Técnicas voluntarias implementadas}

% \subsubsection{Variantes de Búsqueda Local}

% Además de la BL primer mejor obligatoria (\texttt{LOCALSEARCH}), se han añadido
% las dos variantes voluntarias siguientes:

% \begin{itemize}[noitemsep]
%     \item \textbf{BL Best} (\texttt{LOCALSEARCHBEST}): explora el vecindario completo
%           y selecciona el mejor vecino disponible en cada iteración.
%     \item \textbf{BL Multi} (\texttt{LOCALSEARCHMULTI}): ejecuta varias búsquedas
%           locales consecutivas desde distintas soluciones iniciales hasta agotar el
%           presupuesto de evaluaciones.
% \end{itemize}

% Estas dos variantes permiten aislar el efecto de la intensificación: una busca el
% mejor vecino posible en cada paso y la otra intenta amortiguar los óptimos locales
% reiniciando la búsqueda desde diferentes puntos. Estas son las mejoras que se implementaron en la práctica 1.

% \subsubsection{Mutación Gaussiana en AGG}

% Se implementa una variante \textbf{AGG-Gauss} (\texttt{gaussian\_mutation.cpp}) que
% sustituye la mutación por transferencia por una perturbación continua basada en
% ruido normal, generada mediante Box-Muller y seguida de reparación de factibilidad.
% Con ello se compara una mutación estadística suave con la transferencia fija usada
% por los algoritmos genéticos obligatorios.

% \subsubsection{Encadenamiento de Búsqueda Local en AM-LSCh}

% Se añadió \textbf{AM-LSCh} (\texttt{memetic\_lsch.cpp}), que preserva el estado de la
% búsqueda local por individuo. La idea es no reiniciar siempre el refinamiento desde
% cero, sino acumular información sobre qué soluciones ya han sido tratadas y seguir
% encadenando mejoras sobre candidatos prometedores. Esta variante mantiene el mismo
% marco memético, pero cambia la gestión interna de la explotación local.

% \subsubsection{Evolución Diferencial}

% También se incorporó un algoritmo \textbf{DE} clásico (DE/rand/1/bin) con mutación
% diferencial, cruce binomial y selección uno a uno:

% \begin{itemize}[noitemsep]
%     \item Mutación diferencial: $v_i = x_{r_0} + F \cdot (x_{r_1} - x_{r_2})$.
%     \item Cruce binomial con tasa $CR \in [0,1]$.
%     \item Reemplazo determinista del vector objetivo si el trial mejora su fitness.
% \end{itemize}

% Esta familia es relevante porque introduce una dinámica evolutiva distinta a la de
% AGG y AM, y permite comprobar si su forma de explorar el espacio resulta competitiva
% para representaciones reales del portfolio.

% \subsection{Análisis comparativo de los extras}

% Las Figuras~\ref{fig:ibex-todos}--\ref{fig:sp500-todos} muestran los boxplots con
% los quince algoritmos completos, es decir, los diez obligatorios más los cinco
% voluntarios, divididos en dos partes por legibilidad. En casi todos los gráficos
% aparecen uno o dos outliers aislados, lo que indica que incluso cuando la caja
% central es compacta siguen existiendo ejecuciones puntuales alejadas del resto.

% \subsubsection{BL Best y BL Multi frente a la BL base}

% BL Best no mejora a la BL primer mejor; al contrario, se queda por detrás en los
% tramos más importantes. En IBEX 35 empeora alrededor de un 2\% tanto en fitness
% histórico como en validación. En S\&P 100 y S\&P 500 el retroceso es bastante más
% fuerte, del orden de un 20\%--25\% en \textit{2015-2024 Fitness} y de un 15\%--16\%
% en \textit{2025 Fitness}. El patrón es claro: evaluar todo el vecindario no compensa
% si el problema ya está bien resuelto con una estrategia de primer mejora y una
% reparación sólida.

% BL Multi sí ofrece una mejora más razonable. En IBEX 35 la diferencia es pequeña,
% en torno al 1\%, pero en S\&P 100 y S\&P 500 mejora a la BL base entre un 1\% y un
% 3\%, sobre todo en el fitness histórico. Esto sugiere que reiniciar la búsqueda local
% desde distintos puntos ayuda a explorar mejor el espacio, aunque el beneficio no es
% lo bastante grande como para convertirla en una dominadora clara de la versión base.

% \subsubsection{AGG-Gauss frente a los AG obligatorios}

% AGG-Gauss mejora de forma consistente al cruce aritmético. Frente a AGG-Arit, la
% ganancia ronda entre un 3\% y un 6\% según el mercado y la fase medida. En IBEX 35
% está aproximadamente un 3\% por encima en fitness histórico y algo más de un 1\%
% mejor en validación. En S\&P 100 la ventaja sobre el cruce aritmético es aún más
% clara, mientras que en S\&P 500 mantiene esa mejora aunque con una dispersión algo
% mayor.

% Sin embargo, AGG-Gauss no alcanza a AGG-BLX. El cruce BLX sigue siendo la opción más
% robusta entre las dos que se comparan en el informe principal, y AGG-Gauss queda muy
% cerca pero por debajo. En términos prácticos, la mutación gaussiana aporta una
% perturbación suave interesante, pero no consigue superar la ventaja estructural del
% muestreo expandido de BLX para esta representación real.

% \subsubsection{AM-LSCh frente a los meméticos obligatorios}

% AM-LSCh se mueve prácticamente en la misma banda que AM-Best y AM-Rand. En IBEX 35
% las diferencias son mínimas, y en S\&P 100 y S\&P 500 su comportamiento oscila entre
% pequeñas mejoras y pequeñas pérdidas de menos de un 1\%. En otras palabras, el
% encadenamiento de estados no degrada la calidad, pero tampoco produce una ruptura
% clara respecto a las variantes meméticas ya incluidas en la parte obligatoria.

% Desde el punto de vista del análisis, esto es útil porque confirma que la manera de
% organizar la búsqueda local puede cambiar sin penalizar de forma apreciable la
% calidad final. Además, el coste temporal permanece en el mismo orden de magnitud que
% el de AM-Best, así que la elección entre ambas variantes puede hacerse por criterio de
% implementación más que por rendimiento puro.

% \subsubsection{Evolución Diferencial frente a AGG y AM}

% DE es la variante voluntaria que mejor encaja con mercados grandes. En S\&P 100 y
% S\&P 500 mejora a AGG-BLX en \textit{2015-2024 Fitness} en torno a un 2\%--3\%, y
% además en S\&P 500 consigue también una ligera ventaja en \textit{2025 Fitness}. En
% S\&P 100 el comportamiento es algo más mixto: mejora en entrenamiento pero queda muy
% cerca o ligeramente por detrás en validación.

% Frente a AM-Best, DE está prácticamente empatado en IBEX 35 y S\&P 100, y en S\&P 500
% resulta ligeramente mejor tanto en entrenamiento como en validación. Esto refuerza la
% idea de que la evolución diferencial escala bien cuando aumenta la dimensión del
% problema y que, en los casos grandes, puede competir de tú a tú con los mejores
% meméticos de la práctica.

% \subsection{Distribuciones completas}

% Las figuras con todos los algoritmos permiten ver el panorama completo. La parte
% superior de las distribuciones sigue ocupada por las mejores combinaciones del
% informe principal, a las que se suman DE y AM-LSCh en una posición muy competitiva.
% Los peores comportamientos siguen concentrándose en Random y BL Best, lo que encaja
% con lo visto en el análisis cuantitativo.

% También se aprecia que la estabilidad no depende solo del valor medio del fitness.
% Hay algoritmos con medianas atractivas pero con uno o dos outliers marcados, y otros
% que mantienen una caja algo más ancha pero una evolución más homogénea. Por eso en
% esta sección se valora conjuntamente la posición central, la dispersión y la presencia
% de valores atípicos.

% \begin{figure}[H]
%     \centering
%     \begin{minipage}[t]{0.9\textwidth}
%         \centering
%         \includegraphics[width=\textwidth]{figuras/fig_ibex_todos_pt1.png}

%         \smallskip
%         {\small Baselines + AG}
%     \end{minipage}

%     \vspace{0.5cm}

%     \begin{minipage}[t]{0.9\textwidth}
%         \centering
%         \includegraphics[width=\textwidth]{figuras/fig_ibex_todos_pt2.png}

%         \smallskip
%         {\small Meméticos + Extras}
%     \end{minipage}
%     \caption{Distribución del \textit{fitness} con todos los algoritmos en IBEX 35.}
%     \label{fig:ibex-todos}
% \end{figure}

% \begin{figure}[H]
%     \centering
%     \begin{minipage}[t]{0.9\textwidth}
%         \centering
%         \includegraphics[width=\textwidth]{figuras/fig_sp100_todos_pt1.png}

%         \smallskip
%         {\small Baselines + AG}
%     \end{minipage}

%     \vspace{0.5cm}

%     \begin{minipage}[t]{0.9\textwidth}
%         \centering
%         \includegraphics[width=\textwidth]{figuras/fig_sp100_todos_pt2.png}

%         \smallskip
%         {\small Meméticos + Extras}
%     \end{minipage}
%     \caption{Distribución del \textit{fitness} con todos los algoritmos en S\&P 100.}
%     \label{fig:sp100-todos}
% \end{figure}

% \begin{figure}[H]
%     \centering
%     \begin{minipage}[t]{0.9\textwidth}
%         \centering
%         \includegraphics[width=\textwidth]{figuras/fig_sp500_todos_pt1.png}

%         \smallskip
%         {\small Baselines + AG}
%     \end{minipage}

%     \vspace{0.5cm}

%     \begin{minipage}[t]{0.9\textwidth}
%         \centering
%         \includegraphics[width=\textwidth]{figuras/fig_sp500_todos_pt2.png}

%         \smallskip
%         {\small Meméticos + Extras}
%     \end{minipage}
%     \caption{Distribución del \textit{fitness} con todos los algoritmos en S\&P 500.}
%     \label{fig:sp500-todos}
% \end{figure}

% \subsection{Conclusiones sobre las mejoras voluntarias}

% La lectura conjunta de los resultados es bastante coherente con la teoría y con la
% guía de la práctica. La BL base heredada de la práctica 1 sigue siendo una referencia
% muy sólida; BL Best no compensa su mayor intensidad de búsqueda; BL Multi mejora un
% poco la exploración; AGG-Gauss ofrece una mejora clara sobre el cruce aritmético pero
% no supera BLX; AM-LSCh mantiene la calidad de los meméticos sin penalización relevante;
% y DE es la variante que mejor escala cuando el mercado crece, especialmente en
% S\&P 500.

% En resumen, los algoritmos voluntarios no cambian la jerarquía general establecida por
% los obligatorios, pero sí matizan el análisis. Los mejores extras se sitúan muy cerca
% de los mejores obligatorios, mientras que las variantes más agresivas o más exploratorias
% suelen pagar ese cambio con más dispersión o peor validación. Eso encaja con el objetivo
% de esta sección: mostrar que pequeñas modificaciones del operador o de la estrategia de
% explotación pueden mover el rendimiento, pero no siempre en la dirección esperada.
